\section{Aging, Damage, and Imperfection}
\label{sec:aging}

\subsection{Design Stance}

This stage is the one most likely to make the product look cheap, and it is
therefore governed by a stricter rule than the rest of the pipeline: \emph{every
artifact must be off by default, must be physically parameterized, and must be
capable of subtlety}. The failure mode of consumer film apps is not that they
have light leaks; it is that their light leaks are a fixed PNG at fixed opacity
in a fixed corner.

Two structural decisions follow. First, all artifacts are procedural and seeded,
so they never repeat across frames. Second, all artifacts are applied in the
domain where the corresponding physical damage occurs, which for most of them is
the negative, not the print --- meaning a scratch in the negative prints
\emph{light}, not dark, and dust on the negative prints as white specks. Getting
this polarity right is a small detail that separates a considered
implementation from a texture overlay, since almost every consumer application
renders dust as dark.

\subsection{Artifact Taxonomy and Domain}

\begin{table}[htbp]
\caption{Aging Artifacts, Domain, and Polarity}
\label{tab:aging}
\begin{center}
\renewcommand{\arraystretch}{1.15}
\begin{tabularx}{\columnwidth}{@{}l l l >{\raggedright\arraybackslash}X@{}}
\toprule
\textbf{Artifact} & \textbf{Domain} & \textbf{Polarity} & \textbf{Mechanism} \\
\midrule
Dust           & neg. $D$   & $+D$ & Opaque particle blocks printing light; prints white. \\
Scratch (emulsion) & neg. $D$ & $-D$ & Emulsion removed; prints dark. \\
Scratch (base) & neg. $D$   & $+D$ & Base abraded, scatters; prints light. \\
Chemical streak & neg. $D$  & $\pm D$ & Uneven development; direction-aligned. \\
Edge fog       & lin. $E$   & $+E$ & Light leaked at the edges pre-development. \\
Light leak     & lin. $E$   & $+E$ & Broad exposure through a body seam. \\
Halide fog     & neg. $\Dmin$ & $+\Dmin$ & Age/heat/radiation; raises base fog. \\
Expired shift  & $\theta$   & --- & Differential layer speed and gamma loss. \\
Frame jitter   & geometry   & --- & Registration error in transport. \\
Gate weave     & geometry   & --- & Low-frequency lateral drift (video only). \\
\bottomrule
\end{tabularx}
\end{center}
\end{table}

\subsection{Expired Film}

The most rewarding of these to model correctly, because it is entirely a
parameter modulation and costs nothing at render time.

Aging degrades an emulsion in three ways: overall speed loss, differential speed
loss between layers (blue-sensitive layers degrade fastest), and a rise in base
fog. All three are captured by modulating the parameters already defined:

\begin{align}
x_{0,c} &\leftarrow x_{0,c} + \varsigma_a\, \left(0.30 + \upsilon_c\right), \\
\Dmin{}_{,c} &\leftarrow \Dmin{}_{,c} + \varsigma_a\, \phi_a\, \psi_c, \\
\gamma_c &\leftarrow \gamma_c \left(1 - 0.18\,\varsigma_a\right),
\label{eq:expired}
\end{align}

where $\varsigma_a \in [0, 1]$ is the age parameter (interpreted as
"decades past expiry under uncontrolled storage"), $\boldsymbol{\upsilon} =
(0.00, 0.06, 0.19)$ is the differential speed loss, and $\boldsymbol{\psi} =
(0.4, 0.7, 1.0)$ the differential fogging. The result is the familiar look of
expired color negative --- muddy, cyan-shifted shadows, reduced contrast, and a
loss of speed requiring overexposure --- produced entirely by the existing
model. There is no "expired film LUT."

\subsection{Dust and Scratches}

Dust is generated as a Poisson point process with intensity
$\lambda_{\mathrm{dust}}$ particles per square millimetre of film, each particle
given a random radius from a log-normal distribution
$\ln r \sim \mathcal{N}(\mu_r, \sigma_r^2)$ with $\mu_r$ corresponding to a
median particle of $18\,\mu$m. The particle is rendered as a soft-edged disc
adding density:

\begin{equation}
\Delta D_{\mathrm{dust}}(\mathbf{p}) = D_{\mathrm{op}} \sum_i \mathrm{smoothstep}\!\left(r_i, r_i(1-\epsilon), \|\mathbf{p} - \mathbf{p}_i\|\right),
\end{equation}

applied to all three records equally (dust is neutral) and clamped so a particle
cannot exceed the printing system's black.

Because the process is defined in film-plane coordinates with the same
lattice-hashing scheme as the grain (Section~\ref{sec:grain}), dust is
resolution-independent and deterministic under a seed --- and its size relative
to the frame follows the format selection automatically. Dust on 8\,mm is
enormous; dust on 4$\times$5 is invisible. This is correct and free.

Scratches are line segments with a random start point, a strongly
vertically-biased orientation (matching film transport direction), a length drawn
from an exponential distribution, and a width of one to three grains. Each
scratch is rendered with a small perpendicular density profile that is negative
at the core (emulsion removed) and slightly positive at the shoulders (displaced
emulsion), which is what a real scratch looks like under a loupe.

\subsection{Light Leaks and Edge Fog}

Both are applied in the linear exposure domain, which matters: a leak adds
\emph{light}, so it is compressed by the film's shoulder in bright areas and
dominates in dark areas. An overlay applied at the end is uniformly visible and
looks pasted on.

Edge fog is a distance-from-edge falloff,

\begin{equation}
E' = E + \Phi_e \exp\!\left(-\frac{d(\mathbf{p})}{\ell_e}\right)\boldsymbol{\chi}_e,
\end{equation}

with $d$ the distance to the nearest frame edge, $\ell_e$ the penetration depth
in millimetres of film, and $\boldsymbol{\chi}_e$ the leak color (warm, since
light entering a camera body is usually filtered by the body's interior).

Light leaks use a procedural model rather than textures: a leak is parameterized
by an entry point on the frame boundary, a direction, an angular spread, and a
falloff, and is rendered as an anisotropic exponential wedge. Three to five such
wedges with randomized parameters, seeded from the frame seed, produce
non-repeating and physically plausible leaks. A "seam" preset places the entry
points along a line, reproducing the look of a failed back-door light seal, which
is the most common real cause.

\subsection{Chemical Streaking}

Uneven development produces streaks aligned with the direction of development
flow. Model: a one-dimensional noise field $\varpi(y)$ along the transport axis,
smoothed with a long kernel, modulating development activity $A$ from
\eqref{eq:activity}:

\begin{equation}
A(\mathbf{p}) = A \cdot \left(1 + \varsigma_{\mathrm{streak}}\, \varpi(y)\right).
\end{equation}

Because it modulates $A$ rather than density directly, streaking correctly
produces \emph{contrast} variation as well as density variation --- streaks are
not merely lighter or darker bands, they have different gamma. This is the
detail that makes it read as chemistry rather than as a gradient overlay.
Evaluating a spatially varying $A$ requires the parameter modulation of
Section~\ref{sec:development} to be evaluated per-pixel rather than on the host,
so this effect (and only this effect) forces a shader-side evaluation of
\eqref{eq:gammatime} and \eqref{eq:fog}. It is consequently gated behind an
"enable" that recompiles the pipeline variant; see
Section~\ref{sec:specialization}.

\subsection{Frame Geometry Artifacts}

Frame jitter is a per-frame random translation and rotation applied at the
geometric warp stage, with magnitude specified in micrometres of film so that it
scales with format. For still photography a single frame's jitter is
indistinguishable from a crop offset and is therefore disabled by default; it
exists in v1.0 only so that the parameter and its persistence are in place for
the video release, where gate weave --- a low-frequency correlated drift modelled
as a two-octave value-noise walk --- becomes the visible and important effect.

\subsection{Ordering Within the Stage}

Within the aging stage, artifacts are applied in the order:

\begin{enumerate}
\item Parameter-domain (expired shift, halide fog) --- these modify $\theta$ and
are consumed by earlier stages, so they are technically evaluated first despite
being authored here.
\item Exposure-domain (edge fog, light leaks) --- before the characteristic
curve.
\item Density-domain (dust, scratches, streaking) --- after the characteristic
curve, before printing.
\item Geometric (jitter) --- fused into the optical warp.
\end{enumerate}

This ordering means the "Aging" panel in the UI is not a pipeline stage in the
render-graph sense; it is a control panel whose parameters are injected at four
different points. Section~\ref{sec:rendergraph} explains how the graph handles
this, and it is the principal reason the render graph is a graph rather than an
array.
